% =============================================================================
% Diapositivas de respaldo (después de \appendix): no se numeran en el total
% =============================================================================

\begin{frame}{Respaldo: propiedad de convolución}
  \begin{teorema}
    Si $x(t) \fourier X(\j\omega)$ y $h(t) \fourier H(\j\omega)$, entonces
    $x(t) \conv h(t) \fourier X(\j\omega)\, H(\j\omega)$.
  \end{teorema}
  \begin{demostracion}
    \vspace{-\baselineskip}
    \begin{align*}
      \TF{x \conv h}
        &= \int_{-\infty}^{\infty} \left[ \int_{-\infty}^{\infty} x(\tau)\, h(t - \tau) \dd\tau \right]
           \e^{-\j\omega t} \dd t \\
        &= \int_{-\infty}^{\infty} x(\tau) \left[ \int_{-\infty}^{\infty} h(t - \tau)\,
           \e^{-\j\omega t} \dd t \right] \dd\tau
         = \int_{-\infty}^{\infty} x(\tau)\, \e^{-\j\omega\tau} H(\j\omega) \dd\tau \\
        &= X(\j\omega)\, H(\j\omega) \qedhere
    \end{align*}
  \end{demostracion}
\end{frame}
